Dentro de \texttt{lib/}, o código foi dividido em módulos para separar a interface gráfica, as regras de navegação e predição de texto, o registro de métricas, a síntese de voz e o rastreamento ocular. A partir do ponto de entrada \texttt{main.dart}, que inicializa o aplicativo \texttt{ElaCommunicatorApp} e abre a tela principal, a arquitetura organiza-se na camada de telas (\texttt{screens/}) e em cinco pacotes internos (\texttt{widgets/}, \texttt{models/}, \texttt{metrics/}, \texttt{tts/} e \texttt{gaze/}), representados na \refFig{fig:modulos}.

\begin{figure}[htbp]
    \centering
    \includegraphics[width=\textwidth]{modulos}
    \caption{Diagrama de módulos do diretório \texttt{lib/} e suas dependências com pacotes e recursos externos.}
    \label{fig:modulos}
\end{figure}

Como mostra a \refFig{fig:modulos}, a tela principal (\texttt{CommunicatorScreen}) concentra a coordenação entre os módulos, enquanto o acesso à câmera, a inferência dos modelos e a gravação de arquivos permanecem encapsulados em seus respectivos pacotes. As seções a seguir descrevem a responsabilidade de cada módulo e das classes que o compõem.

\section{Ponto de Entrada e Telas (\texttt{main.dart} e \texttt{screens/})}

O ponto de entrada e as duas telas da aplicação controlam o estado geral da sessão e integram os demais módulos:

\begin{itemize}
    \item \texttt{ElaCommunicatorApp}: \textit{widget} raiz definido em \texttt{main.dart}, que configura o componente \texttt{MaterialApp} com tema escuro de alto contraste (Material~3) e abre a \texttt{CommunicatorScreen} como tela inicial.
    \item \texttt{CommunicatorScreen}: tela principal do comunicador (\refFig{fig:tela_teclado}). Mantém o texto da mensagem em composição, a linha e a coluna atualmente selecionadas no teclado, o índice da sugestão destacada na barra superior e o estado do rastreador ocular e do sintetizador de voz. Toda ação de navegação ou seleção chega ao método \texttt{\_handleCommand} na forma de um \texttt{GazeCommand}, acompanhado da origem do comando (olhar ou toque). Para facilitar os testes de interface, a tela recebe opcionalmente por injeção de dependência instâncias de \texttt{SpeechSynthesizer}, \texttt{Word\allowbreak Predictor} e \texttt{Typing\allowbreak Metrics\allowbreak Recorder}, criando as implementações padrão quando executada normalmente no aparelho.
    \item \texttt{CalibrationScreen}: tela aberta a partir do botão de status do rastreador na \texttt{CommunicatorScreen} para calibrar o olhar do usuário (\refFig{fig:tela_calibragem}). Apresenta uma sequência de alvos na tela, recolhe os vetores pupila--canto calculados pelo \texttt{GazeTracker} em cada posição, ajusta os coeficientes da regressão polinomial com descarte de pontos discrepantes e devolve à tela principal o mapeamento resultante (\texttt{GazeMap}).
\end{itemize}

\begin{figure}[!htbp]
    \centering
    \includegraphics[width=0.78\textwidth]{tela_teclado.jpeg}
    \caption{Tela principal do comunicador (\texttt{CommunicatorScreen}) no arranjo de bordas, com a barra de mensagem e sugestões, o botão de status do rastreador, a prévia da câmera frontal, o teclado virtual QWERTY, o botão de confirmação, as setas direcionais nas bordas e o indicador circular do ponto estimado do olhar.}
    \label{fig:tela_teclado}
\end{figure}

\begin{figure}[!htbp]
    \centering
    \includegraphics[width=0.78\textwidth]{tela_calibragem.jpeg}
    \caption{Tela de calibração (\texttt{CalibrationScreen}) durante a coleta de amostras, exibindo o alvo ativo com anel de progresso, a contagem de alvos na parte inferior e o círculo ciano correspondente à estimativa atual do olhar.}
    \label{fig:tela_calibragem}
\end{figure}

  \section{Componentes da Interface Visual (\texttt{widgets/})}

O pacote \texttt{widgets/} contém os componentes visuais da interface, que recebem o estado atual por parâmetros e notificam eventos à tela principal:

\begin{itemize}
    \item \texttt{OnScreenKeyboard}: desenha o teclado virtual QWERTY com uma única tecla destacada por vez, correspondente à linha e à coluna selecionadas, e atualiza o rótulo e a disponibilidade da tecla FALAR conforme o estado da síntese de voz.
    \item \texttt{MessageBar}: faixa superior que exibe a mensagem em composição e, quando há texto digitado, mostra logo abaixo a lista horizontal de sugestões de palavras, destacando aquela que estiver selecionada pelo foco de navegação.
    \item \texttt{EdgeControls}: dispõe as quatro setas direcionais (cima, baixo, esquerda e direita) ao longo das bordas da tela, emoldurando a área central do teclado para maximizar a separação angular entre os alvos de olhar.
    \item \texttt{DirectionPad}: dispõe as quatro setas e o botão de confirmação em uma grade compacta de $3 \times 3$ em formato de cruz, posicionada ao lado ou acima do teclado.
    \item \texttt{GazeButton}: botão compartilhado pelos arranjos \texttt{EdgeControls} e \texttt{DirectionPad}, com estilo visual diferenciado para setas de direção e para a ação de confirmação e realce visual quando o olhar aponta para sua direção. O uso do mesmo componente nos dois arranjos garante que a comparação entre eles avalie apenas a diferença geométrica.
    \item \texttt{CameraView}: \textit{widget} com estado (\texttt{StatefulWidget}) responsável por inicializar a câmera frontal, exibir uma prévia compacta no cabeçalho para auxiliar o posicionamento do rosto e fornecer o \texttt{CameraController} ao rastreador ocular por meio do \textit{callback} \texttt{onControllerReady}. No método \texttt{\_start}, enumera as câmeras do dispositivo (\texttt{availableCameras}), seleciona a câmera frontal (\texttt{CameraLensDirection.\allowbreak front}, recaindo na primeira disponível caso não exista câmera frontal) e instancia o \texttt{CameraController} com três escolhas voltadas ao pipeline de rastreamento: resolução média de 480p (\texttt{ResolutionPreset.\allowbreak medium}), pois o custo de inferência do ML Kit Face Mesh cresce com a resolução e domina o tempo de processamento (cerca de 150~ms por quadro a 720p contra 48~ms da detecção dos dois olhos juntos, enquanto a 480p cada olho ainda ocupa cerca de 100~pixels de largura, acima dos aproximadamente 70~pixels das imagens de treinamento); áudio desativado (\texttt{enableAudio: false}), evitando solicitar permissão de microfone desnecessária; e formato de imagem \texttt{ImageFormatGroup.\allowbreak nv21}, que entrega os quadros brutos exatamente na codificação exigida pelo ML Kit no Android e amostrada por \texttt{Nv21Frame}, dispensando conversões de cor do quadro inteiro. A classe de estado incorpora o \textit{mixin} \texttt{WidgetsBindingObserver} para acompanhar o ciclo de vida do aplicativo (\texttt{didChangeAppLifecycleState}): ao transitar para \texttt{AppLifecycleState.\allowbreak inactive} (como ao ir para segundo plano ou bloquear a tela), invoca \texttt{\_release}, que notifica \texttt{onControllerReady(null)} para que a \texttt{CommunicatorScreen} encerre o \texttt{GazeTracker} e descarta o controlador antes que o Android derrube a sessão da câmera e congele a prévia; ao retornar a \texttt{AppLifecycleState.\allowbreak resumed}, chama \texttt{\_start} para reabrir a câmera automaticamente (protegido pela trava \texttt{\_initializing} contra chamadas concorrentes e por verificações de \texttt{mounted}). Na renderização, exibe um indicador de progresso enquanto o controlador inicializa, um estado de erro com a descrição da \texttt{CameraException} e botão de nova tentativa caso a abertura falhe (por exemplo, se a permissão for negada) ou o \texttt{CameraPreview} envolvido em um \texttt{FittedBox} com \texttt{BoxFit.cover}, que preenche a moldura de proporção $3:4$ do cabeçalho recortando as bordas excedentes sem distorcer o rosto do usuário.
\end{itemize}

  \section{Modelos e Regras da Interface (\texttt{models/})}

O pacote \texttt{models/} reúne as definições de dados, a geometria de navegação do teclado e o modelo de preenchimento automático, implementados em Dart puro sem dependência de pacotes nativos:

\begin{itemize}
    \item \texttt{GazeCommand}
    enumeração com os cinco comandos que controlam a interface (\texttt{up}, \texttt{down}, \texttt{left}, \texttt{right} e \texttt{select}). Funciona como contrato único de entrada: os botões na tela, o rastreador ocular e o digitador simulado produzem os mesmos comandos, de modo que a interface não depende da forma como o comando foi gerado.
    \item \texttt{ControlStyle}
    enumeração que seleciona entre o arranjo de setas nas bordas da tela (\texttt{edges}) e o arranjo em cruz (\texttt{cross}).
    \item \texttt{KeyboardLayout} e \texttt{KeyDef}
    definem a grade do teclado (\texttt{kKeyboardRows}), composta pelas três linhas das 26 letras do padrão QWERTY e por uma quarta linha com as quatro teclas de ação (FALAR, ESPAÇO, APAGAR e LIMPAR), dimensionadas por pesos proporcionais (\texttt{flex}).
    \item \texttt{KeyboardNavigation}
    implementa as regras de movimentação do foco sobre a grade do teclado, limitando o índice da coluna ao comprimento da linha de destino quando o foco muda de linha. Fornece a localização de qualquer tecla (\texttt{find}), o conjunto de teclas vizinhas a um comando de distância (\texttt{neighbors}) e o caminho mínimo de comandos entre duas teclas por busca em largura (\texttt{path}, Seção~\ref{sec:bfs}), sendo compartilhada pela tolerância a erros de digitação, pelo digitador simulado e pelo cálculo da linha de base nas métricas (Equação~\ref{eq:custo_base}).
    \item \texttt{WordPredictor}
    motor de preenchimento automático de palavras baseado em trigramas com suavização \textit{Stupid Backoff} e tolerância a erros de digitação. Carrega o \textit{corpus} de frases de \texttt{assets/phrases/}, sugere até três palavras que completam o prefixo em digitação ou antecipam a próxima palavra (\texttt{suggest} e \texttt{predict}) e insere a sugestão escolhida na mensagem (\texttt{accept}). Seu funcionamento completo é detalhado na Seção~\ref{sec:met_preenchimento}.
\end{itemize}
\section{Preenchimento automático de palavras} \label{sec:met_preenchimento}
Presente em \texttt{models/word\_predictor.dart}, o preenchimento automático foi implementado como um módulo independente da interface, a classe \texttt{WordPredictor}, escrita em Dart sem dependência do Flutter. Isso permite que o mesmo código seja executado tanto no aplicativo quanto em ferramentas de linha de comando, como a usada na avaliação offline. O módulo recebe o texto digitado até o momento e devolve as três palavras mais prováveis segundo um modelo de trigramas com \textit{Stupid Backoff} (Seção~\ref{sec:suavizacao}), treinado em um \textit{corpus} de frases de \gls{CAA}. As subseções a seguir descrevem o \textit{corpus} (Seção~\ref{sec:met_corpus}), o tratamento do texto (Seção~\ref{sec:met_normalizacao}), a construção do modelo (Seção~\ref{sec:met_modelo}), a geração e a inserção das sugestões (Seção~\ref{sec:met_sugestoes}), a tolerância a erros de digitação (Seção~\ref{sec:met_tolerancia}) e o método de avaliação offline (Seção~\ref{sec:met_avaliacao}). Os resultados da avaliação são apresentados no \refCap{5_Resultados}.

\subsection{\textit{Corpus}} \label{sec:met_corpus}

Como um modelo de n-gramas reflete o gênero dos textos em que foi treinado (Seção~\ref{sec:limitacoes_ngramas}), o \textit{corpus} é formado apenas por frases da comunicação cotidiana de uma pessoa com limitação motora grave. São 565 frases, de duas origens:
\begin{itemize}
    \item 15 frases curtas escritas no início do desenvolvimento e mantidas no próprio código, como ``eu quero beber água'' e ``preciso de ajuda'';
    \item 550 frases distribuídas em cinco arquivos JSON, incluídos no aplicativo como recursos (\textit{assets}), o que permite ampliar o \textit{corpus} sem alterar o código.
\end{itemize}

As 550 frases dos arquivos foram geradas com o auxílio de uma ferramenta de inteligência artificial generativa. % TODO: nome da ferramenta e se as frases foram revisadas manualmente
Elas cobrem sintomas e desconfortos (``Estou com dor de cabeça'', ``Estou com falta de ar''), cuidados e posicionamento do corpo (``Mova o meu braço direito'', ``Ajuste a máscara de oxigênio''), controle do ambiente (``Abra a janela''), entretenimento (``Coloque um áudio para eu ouvir''), contato com outras pessoas (``Ligue para a minha sobrinha''), respostas curtas (``Sim'', ``Não'', ``Talvez'') e o controle da própria comunicação (``Terminei a frase'', ``Esquece o que eu disse''). No total, o \textit{corpus} soma 2.695 palavras, com média de 4,8 palavras por frase e vocabulário de 767 palavras distintas, e não contém frases repetidas.

O uso de frases geradas é uma limitação deste trabalho: elas não foram produzidas por usuários de \gls{CAA}, e o vocabulário e a frequência das palavras refletem as escolhas do gerador, e não o uso real. Por isso, a avaliação offline (Seção~\ref{sec:met_avaliacao}) mede a capacidade do modelo de prever frases desse mesmo conjunto, e não necessariamente a comunicação de pessoas com \gls{ELA}, que só pode ser observada nos testes com usuários (\refCap{5_Resultados}).

\subsection{Tratamento do Texto} \label{sec:met_normalizacao}

O teclado do aplicativo tem apenas as 26 letras, sem acentos, e quatro teclas de ação (FALAR, ESPAÇO, APAGAR e LIMPAR), enquanto o \textit{corpus} é escrito com acentos e pontuação. Para que o texto digitado corresponda às palavras do \textit{corpus}, ambos passam pela mesma normalização: o texto é dividido em palavras nos sinais de pontuação, as palavras são convertidas para minúsculas, e os acentos, a cedilha e os hífens são removidos. Assim, ``ÁGUA'' e ``água'' tornam-se ``agua'', e ``ligue o ar-condicionado.'' torna-se a sequência ``ligue'', ``o'', ``arcondicionado''. As palavras com hífen são mantidas inteiras porque, sem uma tecla de hífen, o usuário não tem como digitar suas partes separadamente.

A normalização não se perde na exibição: para cada palavra normalizada, o modelo guarda a grafia original mais frequente no \textit{corpus}. A sugestão é mostrada como ``água'' mesmo que o usuário tenha digitado ``AGU'' e, ao ser aceita, é inserida em maiúsculas, como o restante da mensagem, e com o acento (``ÁGUA'').

\subsection{Modelo de Trigramas} \label{sec:met_modelo}

O modelo é construído na inicialização do aplicativo a partir das frases do \textit{corpus}, sem arquivo de modelo pré-treinado. Cada frase normalizada recebe dois símbolos \texttt{<s>} no início (Seção~\ref{sec:ngramas}) e, para cada palavra, são incrementadas três tabelas de contagem: a da palavra após as duas anteriores (trigrama), após a anterior (bigrama) e sem contexto (unigrama), junto com o total de cada contexto. Essas tabelas fornecem diretamente os numeradores e denominadores da Equação~\ref{eq:mle_trigrama}. Diferentemente do usual~\cite[cap.~3]{jurafsky2026slp}, não é acrescentado o símbolo de fim de frase \texttt{</s>}, pois o fim da mensagem é indicado pelo próprio usuário, com a tecla FALAR, e nunca precisa ser sugerido.

As candidatas são pontuadas pelo \textit{Stupid Backoff} (Equação~\ref{eq:stupid_backoff}) com $\alpha = 0{,}4$, o valor usado por Brants \textit{et al.}~\cite{brants2007large}. O modelo consulta primeiro o contexto de duas palavras e atribui a cada palavra observada sua frequência relativa; em seguida, consulta o contexto de uma palavra, pontuando apenas as palavras ainda não pontuadas e multiplicando o resultado por $\alpha$; por fim, recorre aos unigramas, com fator $\alpha^2$. Empates são desfeitos pela ordem alfabética, o que torna as sugestões determinísticas. Com um vocabulário de 767 palavras, as pontuações são recalculadas a cada tecla, sem armazenamento intermediário.

A escolha de um modelo de n-gramas, e não de um modelo neural, segue a discussão da Seção~\ref{sec:limitacoes_ngramas}: o modelo é construído em memória a partir de um \textit{corpus} pequeno, funciona sem conexão com a internet e é executado no próprio \textit{tablet}. A ordem do modelo e o fator $\alpha$ são parâmetros do módulo e podem ser alterados sem mudanças no restante do código.

\subsection{Trie (trigramas)}

\subsection{Geração e Inserção das Sugestões} \label{sec:met_sugestoes}

As três sugestões ficam em uma barra acima do teclado. A partir da primeira linha do teclado, o comando ``cima'' leva o foco à barra, os comandos ``esquerda'' e ``direita'' percorrem as sugestões e o comando de seleção insere a sugestão destacada. Cada sugestão a mais na lista exigiria mais comandos para ser alcançada, o que justifica uma lista curta (Seção~\ref{sec:preenchimento}). A forma de gerar as sugestões depende do texto digitado:
\begin{itemize}
    \item \textbf{Mensagem vazia ou terminada em espaço:} o modelo sugere a próxima palavra, usando como contexto as duas últimas palavras (ou os símbolos \texttt{<s>}). Com a mensagem vazia, por exemplo, as sugestões são ``estou'', ``coloque'' e ``me'', as palavras que mais iniciam frases no \textit{corpus}.
    \item \textbf{Mensagem terminada no meio de uma palavra:} o modelo completa a palavra, considerando apenas as candidatas que começam pelas letras digitadas e pontuando-as com o contexto das palavras anteriores. Após ``ESTOU COM D'', as sugestões são ``dor'', ``dificuldade'' e ``de''.
    \item \textbf{Caso ambíguo:} sem o espaço final, a última palavra digitada pode estar incompleta ou já ser uma palavra inteira. Após ``E'', o usuário pode estar começando ``estou'' ou ter digitado a conjunção ``e''.
\end{itemize}

No caso ambíguo, as duas leituras competem pela pontuação da sequência inteira. Sendo $h$ o histórico de palavras anteriores, $u$ a última palavra digitada, $c$ uma palavra que a completa e $n$ uma possível palavra seguinte, as pontuações são
\begin{equation}
    S_{\textrm{completar}}(c) = S(c \mid h) \qquad \textrm{e} \qquad S_{\textrm{próxima}}(n) = S(u \mid h)\,S(n \mid h\,u),
    \label{eq:ambiguidade}
\end{equation}
e as três maiores, de qualquer das duas leituras, são exibidas. A leitura de $u$ como palavra inteira só é considerada se o contexto formado por $u$ e a palavra anterior (ou o símbolo \texttt{<s>}) tiver sido observado no \textit{corpus}. Após ``E'', as sugestões são ``estou'' e ``eu'', que completam a letra digitada, e ``a'', que viria depois da conjunção ``e''.

Cada sugestão é devolvida com a indicação de qual das leituras a gerou, e uma única função do módulo, \texttt{accept}, a insere no texto: se a sugestão completa a palavra, ela substitui a parte digitada; caso contrário, é acrescentada depois dela. Em ambos os casos, a palavra é inserida em maiúsculas e seguida de um espaço, de modo que as próximas sugestões já sejam de próxima palavra. Assim, aceitar ``beber'' após ``EU QUERO B'' ou após ``EU QUERO'' produz o mesmo texto, ``EU QUERO BEBER''. Essa indicação precisa vir do modelo, porque o texto sozinho não distingue os casos: após ``E'', aceitar ``estou'' substitui a letra e produz ``ESTOU'', enquanto aceitar ``a'' a mantém e produz ``E A''. A mesma função é usada pela tela do aplicativo e pelo digitador simulado da avaliação offline.

\subsection{Tolerância a Erros de Digitação} \label{sec:met_tolerancia}

Na forma descrita até aqui, a completação exige que as letras digitadas sejam o início exato da palavra. Um único erro anula as sugestões: após ``EU QUERO BW'', digitado no lugar de ``EU QUERO BE'', nenhuma palavra do \textit{corpus} é compatível. Para corrigir o erro, o usuário precisa levar o foco até a tecla APAGAR, na última linha do teclado, e voltar às letras, e cada um desses passos é um comando dado com o olhar.

Como o teclado é percorrido por setas, os erros esperados não são os de um teclado físico. O mais provável é a troca por uma tecla vizinha na grade, causada por uma seta a mais ou a menos ou por uma seleção disparada um passo antes ou depois. Também são possíveis a letra repetida, quando a seleção dispara duas vezes, e a letra omitida. A troca de posição entre duas letras, comum em teclados físicos, exigiria navegar de uma letra à outra e foi desconsiderada.

Para comparar as letras digitadas com as palavras do vocabulário, foi usada a distância de edição mínima, definida como o menor número de operações de inserção, remoção e substituição necessárias para transformar uma sequência de caracteres em outra; na distância de Levenshtein, cada operação tem custo 1~\cite[cap.~2]{jurafsky2026slp}. A distância é calculada por programação dinâmica: sendo $D[i,j]$ a distância entre os $i$ primeiros caracteres da sequência de origem e os $j$ primeiros da de destino, cada célula da tabela é obtida a partir das três vizinhas já calculadas, e o algoritmo admite pesos arbitrários para as operações. Em corretores ortográficos, por exemplo, substituições entre letras vizinhas no teclado são mais prováveis~\cite[cap.~2]{jurafsky2026slp}.

Neste trabalho, o método foi adaptado em dois pontos. Primeiro, como o usuário ainda está digitando, as letras digitadas $x$ não são comparadas com a palavra candidata $w$ inteira, mas com o seu início mais próximo:
\begin{equation}
    d(x, w) = \min_{0 \le j \le |w|} D(x,\, w_{1:j}),
    \label{eq:distancia_prefixo}
\end{equation}
o que não exige cálculo adicional, pois a última linha da tabela de programação dinâmica já contém a distância até cada início $w_{1:j}$. Com $d = 0$, recai-se no caso exato. Segundo, os custos refletem o teclado: a substituição por uma tecla vizinha, a uma seta de distância na grade, custa 1, assim como a inserção e a remoção de uma letra, e a substituição por uma tecla distante custa 2. As vizinhas são obtidas pela mesma regra de navegação da tela; as da letra E, por exemplo, são W, R e D.

São toleradas palavras com $d(x, w) \le 1$, ou seja, com no máximo um erro, e apenas a partir de três letras digitadas, já que, com uma ou duas, quase todas as palavras do vocabulário estariam a um erro de distância. A pontuação de uma candidata com erro é a do \textit{Stupid Backoff} multiplicada por um fator de penalidade $\beta$ elevado ao custo:
\begin{equation}
    S_{\textrm{erro}}(w) = S(w \mid h)\,\beta^{\,d(x, w)}, \qquad \beta = 0{,}01,
    \label{eq:penalidade}
\end{equation}
de modo que uma candidata exata quase sempre fica à frente de uma com erro, e as candidatas com erro ocupam, na prática, as vagas que sobram na lista. Com esses parâmetros, após ``EU QUERO BWB'' o modelo sugere ``beber'', já que W é vizinha de E, e a função \texttt{accept} substitui a palavra parcial inteira, com o erro, produzindo ``EU QUERO BEBER''. O mínimo de três letras e o fator $\beta$ foram escolhidos por uma varredura de valores na avaliação offline, apresentada no \refCap{5_Resultados}.

\subsection{Avaliação Offline} \label{sec:met_avaliacao}

Para medir o preenchimento automático antes dos testes com usuários e comparar versões do modelo, foi construída uma avaliação offline, sem usuários, em que um digitador simulado escreve as frases do \textit{corpus}. O digitador segue as mesmas regras da tela do aplicativo:
\begin{itemize}
    \item o foco começa na primeira tecla, com a mensagem vazia, e a frase termina com a tecla FALAR;
    \item a navegação é feita por setas e, como as linhas do teclado têm comprimentos diferentes, a coluna é limitada ao tamanho da linha de destino;
    \item a barra de sugestões é alcançada com o comando ``cima'' a partir da primeira linha.
\end{itemize}
O digitador percorre sempre o menor caminho de setas entre duas teclas, obtido por busca em largura (Seção~\ref{sec:bfs}), e aceita uma sugestão assim que ela custa, em comandos, no máximo o mesmo que terminar a palavra letra a letra. Os eventos que ele gera passam pelo mesmo módulo de métricas das sessões reais (Seção~\ref{sec:met_metricas}), de modo que as métricas da simulação e dos testes com usuários têm definições idênticas. Como o digitador conhece a frase e, no cenário básico, não erra, os resultados representam um limite superior da economia possível, e não uma previsão do desempenho de uma pessoa. A velocidade em WPM é estimada supondo 1,5~s por comando, o tempo de permanência padrão do aplicativo, e não é uma medida.

Para que o modelo nunca seja avaliado em uma frase que viu no treinamento, foi usada validação cruzada em 10 partes. As frases são deduplicadas pela forma normalizada, embaralhadas com uma semente fixa, para que o resultado seja reprodutível, e divididas em 10 partes; para cada parte, um modelo é treinado com as outras nove e o digitador escreve as frases da parte separada. Também é registrada a fração das palavras de teste ausentes do vocabulário de treino, que nenhuma sugestão pode acertar. Como referência, a mesma simulação é feita nas próprias frases de treino, o que mostra quanto o desempenho depende de o modelo já conhecer a frase.

Para avaliar a tolerância a erros, o digitador tem também um modo com erros, em que cada letra é trocada, com probabilidade $p$, por uma letra vizinha na grade, sorteada com semente fixa. São simulados dois comportamentos: no primeiro, o digitador percebe o erro logo após cometê-lo; no segundo, continua a palavra e só o percebe ao terminá-la. Nos dois casos, a cada passo ele escolhe a opção que custa menos comandos: aceitar a sugestão correta, se ela aparecer, ou apagar até o erro e redigitar. O modo com erros simula apenas trocas por teclas vizinhas; letras repetidas e omitidas são toleradas pelo modelo, mas não simuladas.

A avaliação é executada por uma ferramenta de linha de comando (\texttt{dart run tool/\allowbreak evaluate\_\allowbreak predictor.dart}), que imprime uma tabela com os cenários e pode salvar os resultados em JSON. Além dela, testes automatizados verificam a normalização, o \textit{backoff}, a geração e a inserção das sugestões e a tolerância a erros, e garantem que alterações no código não mudem o comportamento esperado.

\section{Rastreamento Ocular (\texttt{gaze/})} \label{sec:met_eyetracking}

O pacote \texttt{gaze/} implementa o pipeline de estimativa do olhar e emissão de comandos por tempo de permanência fundamentado nas Seções~\ref{sec:eyetracking} a~\ref{sec:kalman}:

\begin{itemize}
    \item \texttt{GazeTracker}: coordena o processamento a cada quadro recebido do fluxo da câmera (\texttt{startImageStream}), descartando novos quadros enquanto o anterior ainda está em execução para evitar atraso acumulado. Aciona periodicamente o detector de malha facial do ML Kit, suaviza a posição dos cantos dos olhos pela mediana das últimas leituras, comanda o recorte e a detecção da pupila nos dois olhos, calcula o vetor pupila--canto médio, aplica a regressão calibrada e o filtro de Kalman e converte o ponto estimado no comando da região correspondente quando o tempo de permanência (\textit{dwell}) se completa. Expõe os fluxos \texttt{samples} (estado quadro a quadro) e \texttt{commands} (comandos confirmados).
    \item \texttt{CropParams} e \texttt{CropTransform}: calculam a transformação afim direta e inversa que alinha horizontalmente os cantos do olho e extrai o recorte quadrado de $128 \times 128$ pixels, utilizando exatamente os mesmos parâmetros geométricos do treinamento lidos de \texttt{crop\_params.json} e convertendo o centro detectado da pupila de volta para as coordenadas do quadro.
    \item \texttt{Nv21Frame}: opera diretamente sobre o \textit{buffer} bruto \texttt{NV21} da câmera (Seção~\ref{sec:nv21}). Converte as coordenadas dos pontos do ML Kit do espaço orientado em pé para o espaço bruto do sensor (\texttt{toRaw}, a partir da rotação calculada por \texttt{rotationDegreesFor}) e amostra os pixels do recorte ocular aplicando a transformação afim inversa e a conversão de cores YUV para RGB normalizado (Seção~\ref{sec:conversao_cores}) diretamente no \textit{buffer} de entrada do detector (\texttt{sampleCrop}), sem converter o quadro inteiro.
    \item \texttt{PupilDetector}: executa o modelo \texttt{YOLO26n} exportado para o formato LiteRT (\texttt{pupil\_\allowbreak 128.tflite}) por meio do pacote \texttt{tflite\_flutter}. Lê e escreve diretamente em \textit{buffers} planos tipados (\texttt{Float32List}) para evitar o custo de conversão de listas em Dart e devolve o centro da caixa delimitadora de maior confiança em cada olho.
    \item \texttt{GazeMap}: calcula o vetor pupila--canto normalizado (Equação~\ref{eq:vetor_pupila_canto}), ajusta os coeficientes do polinômio de segunda ordem (Equação~\ref{eq:regressao_polinomial}) por mínimos quadrados com regularização \textit{ridge} resolvidos por eliminação de Gauss (\texttt{fitGazeMap}, Seções~\ref{sec:minimos_quadrados} e~\ref{sec:gauss}), descarta até dois alvos discrepantes na calibração (\texttt{fitGazeMapRobust}) e traduz o ponto estimado na tela em um \texttt{GazeCommand} conforme as faixas de borda (\texttt{GazeBands} e \texttt{commandFor}).
    \item \texttt{GazeKalman}: aplica dois filtros de Kalman unidimensionais independentes, um para cada eixo da tela, segundo o modelo de velocidade constante da Equação~\ref{eq:kalman_matrizes}, utilizando o intervalo de tempo real $\Delta t$ entre quadros consecutivos para suavizar o ponto antes da verificação de permanência.
\end{itemize}

\begin{figure}[!htbp]
    \centering
    \includegraphics[width=\textwidth]{pipeline}
    \caption{Fluxo de processamento do rastreamento ocular a cada quadro da câmera frontal, dividido nas etapas de detecção visual (azul) e de mapeamento geométrico, filtragem e emissão de comando (verde).}
    \label{fig:pipeline}
\end{figure}

Como sintetiza a \refFig{fig:pipeline}, o processamento transforma cada quadro bruto da câmera frontal em um \texttt{GazeCommand} por meio de duas etapas consecutivas: a localização das referências anatômicas na imagem (cantos dos olhos e centro da pupila) e o mapeamento geométrico e temporal para as regiões da interface. As subseções a seguir detalham a captura dos quadros e a detecção dos cantos oculares (Seção~\ref{sec:met_facemesh}), o recorte alinhado e a amostragem direta do \textit{buffer} \texttt{NV21} (Seção~\ref{sec:met_recorte}), o treinamento e a inferência do detector de pupila (Seção~\ref{sec:met_pupila}), a calibração e a regressão polinomial (Seção~\ref{sec:met_calibracao}) e a filtragem de Kalman aliada à seleção por tempo de permanência (Seção~\ref{sec:met_kalman_dwell}).

\subsection{Captura de Quadros e Detecção dos Cantos dos Olhos} \label{sec:met_facemesh}

Assim que o componente \texttt{CameraView} conclui a inicialização da câmera frontal em 480p, a \texttt{CommunicatorScreen} vincula o \texttt{CameraController} ao método \texttt{attach} do \texttt{GazeTracker}, que inicia o fluxo contínuo de imagens (\texttt{startImageStream}). Como a câmera entrega quadros em uma taxa superior ao tempo de inferência combinada dos modelos na CPU, o método \texttt{\_onFrame} mantém uma trava booleana (\texttt{\_busy}): qualquer quadro recebido enquanto o anterior ainda está em processamento é descartado imediatamente. Sem esse descarte, os quadros se acumulariam em fila e a posição estimada na tela apresentaria um atraso crescente em relação ao movimento ocular real.

Para que o ML Kit Face Mesh (Seção~\ref{sec:mlkit}) localize o rosto corretamente em um \textit{tablet}, é indispensável informar a rotação exata que deixa a imagem em pé. No Android, o sensor físico da câmera possui uma orientação fixa em relação ao aparelho (\texttt{sensorOrientation}, igual a $270^\circ$ no \textit{tablet} utilizado), mas a orientação efetiva da imagem depende também de como o dispositivo está sendo segurado (\texttt{deviceOrientation}, igual a $90^\circ$ no modo paisagem). Informar apenas a orientação do sensor quando o \textit{tablet} opera em paisagem faz com que o ML Kit procure rostos deitados a $90^\circ$, falhando na quase totalidade dos quadros. A função \texttt{rotationDegreesFor} resolve essa composição somando o giro da tela no caso da câmera frontal (devido ao espelhamento horizontal) e subtraindo-o na câmera traseira:
\begin{equation}
    \theta_{\text{rotação}} =
    \begin{cases}
        (\theta_{\text{sensor}} + \theta_{\text{tela}}) \bmod 360^\circ & \text{na câmera frontal,} \\
        (\theta_{\text{sensor}} - \theta_{\text{tela}} + 360^\circ) \bmod 360^\circ & \text{na câmera traseira.}
    \end{cases}
    \label{eq:rotacao_camera}
\end{equation}

Da malha de 468 vértices retornada pelo ML Kit, o rastreador extrai apenas dois pares de pontos (\texttt{kEyeCorners}), ilustrados na \refFig{fig:mlkit_facemesh}: os cantos externo e interno do olho direito (índices $33$ e $133$) e os cantos interno e externo do olho esquerdo (índices $362$ e $263$). A ordem dos índices nos dois olhos foi escolhida de modo que o segmento $\mathbf{c}_A \to \mathbf{c}_B$ aponte sempre no sentido positivo do eixo horizontal da imagem, mantendo a convenção do conjunto MPIIGaze e evitando inversões no cálculo do ângulo de inclinação. Como o ML Kit devolve as coordenadas dos vértices no espaço já rotacionado (em pé), enquanto os pixels permanecem na disposição original do sensor, o método \texttt{\_cornersFrom} aplica a transformação inversa (\texttt{toRaw}) a cada vértice antes de armazená-lo, garantindo que os marcos faciais e a amostragem do recorte operem exatamente no mesmo sistema de coordenadas do \textit{buffer} bruto.

\begin{figure}[!htbp]
    \centering
    \includegraphics[width=0.72\textwidth]{mlkit_facemesh.png}
    \caption{Detecção da malha facial de 468 pontos pelo ML Kit Face Mesh (acima) e ampliação da região ocular (abaixo), destacando os quatro vértices dos cantos dos olhos ($\#33$, $\#133$, $\#362$ e $\#263$), o eixo intercantal $\mathbf{c}_A \to \mathbf{c}_B$, o ponto médio $\mathbf{m}$ e o quadrado do recorte alinhado de lado $1{,}6\,\lVert\mathbf{c}_B - \mathbf{c}_A\rVert$.}
    \label{fig:mlkit_facemesh}
\end{figure}

Duas estratégias adicionais tratam o custo computacional e o ruído desses pontos:
\begin{itemize}
    \item \textbf{Reaproveitamento entre quadros:} a inferência do Face Mesh em um quadro de 480p leva cerca de 100~ms quando há um rosto na cena, enquanto a detecção das duas pupilas juntas consome cerca de 38~ms. Como os cantos das pálpebras acompanham a cabeça e se deslocam lentamente entre quadros consecutivos, ao passo que a pupila realiza movimentos rápidos, o \texttt{GazeTracker} executa o Face Mesh apenas uma vez a cada \texttt{meshEveryNFrames} $= 3$ quadros, reutilizando os cantos nos dois quadros intermediários --- salvo se o último conjunto de cantos tiver mais de \texttt{maxCornerAge} $= 400$~ms ou se a orientação do aparelho mudar. Essa decisão reduz o custo médio do estágio facial para cerca de 28~ms por quadro e eleva a taxa global do pipeline de 4--6~qps para 10--12~qps.
    \item \textbf{Mediana temporal dos cantos:} como o vetor pupila--canto (Equação~\ref{eq:vetor_pupila_canto}) divide o deslocamento da pupila pela distância $\lVert \mathbf{c}_B - \mathbf{c}_A \rVert$, qualquer oscilação nos cantos entra diretamente no sinal. Medições no dispositivo mostraram que, embora o tremor típico do ML Kit fique entre 0,1 e 1,8~pixel, ocorrem picos esporádicos de até 13~pixels --- o que, em um olho de cerca de 100~pixels de largura, equivale a um salto de $0{,}13$ unidades em $\mathbf{v}$, suficiente para corromper um alvo de calibração inteiro. Para eliminar esses picos sem introduzir o atraso de uma média móvel longa, o rastreador mantém as últimas \texttt{cornerHistory} $= 5$ leituras do Face Mesh e calcula, em \texttt{\_smoothedCorners}, a mediana independente das coordenadas $x$ e $y$ de cada um dos quatro cantos.
\end{itemize}

\subsection{Recorte Alinhado e Amostragem Direta do \textit{Buffer} \texttt{NV21}} \label{sec:met_recorte}

A partir dos cantos suavizados $\mathbf{c}_A = (x_A, y_A)$ e $\mathbf{c}_B = (x_B, y_B)$ de cada olho no espaço bruto do sensor, o módulo \texttt{eye\_crop.dart} constrói uma transformação afim que isola a região ocular em uma imagem quadrada de $S \times S$ pixels ($S = 128$), compensando tanto a distância do usuário à câmera quanto a inclinação lateral da cabeça (\textit{roll}), conforme ilustrado pelos quadrados azuis na \refFig{fig:mlkit_facemesh}. Sejam $\mathbf{m} = (\mathbf{c}_A + \mathbf{c}_B)/2 = (c_x, c_y)$ o ponto médio entre os cantos, $w_{\text{olho}} = \lVert \mathbf{c}_B - \mathbf{c}_A \rVert$ a largura aparente do olho em pixels e $\theta = \operatorname{atan2}(y_B - y_A,\, x_B - x_A)$ o ângulo do eixo intercantal. Definindo a razão de recorte $r = 1{,}6$ (\texttt{cropRatio}), o lado do recorte no quadro original mede $r\,w_{\text{olho}}$, o que resulta no fator de escala $s = S / (r\,w_{\text{olho}})$. Com $\alpha = s\cos\theta$ e $\beta = s\sin\theta$, a transformação afim direta $\mathbf{M}$ que leva um ponto $(x, y)$ do quadro bruto para $(u, v)$ no recorte é dada por
\begin{equation}
    \begin{pmatrix} u \\ v \end{pmatrix} =
    \begin{pmatrix} \alpha & \beta \\ -\beta & \alpha \end{pmatrix}
    \begin{pmatrix} x - c_x \\ y - c_y \end{pmatrix} +
    \begin{pmatrix} S/2 \\ S/2 \end{pmatrix},
    \label{eq:afim_recorte}
\end{equation}
que posiciona $\mathbf{m}$ no centro $(64, 64)$ do recorte e alinha horizontalmente os dois cantos, separados por uma distância fixa de $S / r = 80$~pixels independentemente da distância do rosto ao \textit{tablet}. A classe \texttt{CropTransform} armazena tanto os seis coeficientes de $\mathbf{M}$ quanto os de sua inversa analítica $\mathbf{M}^{-1}$, usada para amostrar os pixels da imagem e para converter o centro da pupila de volta às coordenadas do quadro (\texttt{toFrame}). Como qualquer diferença entre o recorte aplicado no aplicativo e o utilizado para gerar os dados de treinamento degradaria a precisão do detector sem produzir erros de execução, os parâmetros ($S = 128$, $r = 1{,}6$, \texttt{pupilBoxRatio} $= 0{,}20$ e \texttt{alignRoll} $= \texttt{true}$) são carregados do arquivo \texttt{assets/models/crop\_params.json} exportado pelo próprio \textit{script} de treinamento, e um teste automatizado (\texttt{eye\_crop\_test.dart}) confere os coeficientes de $\mathbf{M}$ e $\mathbf{M}^{-1}$ contra os valores de referência gerados em Python com \texttt{OpenCV}.

Em vez de converter o quadro completo de $640 \times 480$ pixels (307.200 pixels) de \texttt{NV21} para RGB e depois recortar os olhos, a classe \texttt{Nv21Frame} realiza a amostragem inversa diretamente sobre o \textit{buffer} bruto em \texttt{sampleCrop}, executando apenas $2 \times 128 \times 128 = 32.768$ leituras por quadro e dispensando alocações intermediárias. Para cada pixel $(d_x, d_y) \in \{0, \ldots, 127\}^2$ do recorte, avalia-se $\mathbf{M}^{-1}$ no centro do pixel $(d_x + 0{,}5,\, d_y + 0{,}5)$, arredonda-se o resultado para o inteiro mais próximo (interpolação por vizinho mais próximo) e limitam-se as coordenadas $(s_x, s_y)$ aos limites do quadro bruto $[0, W-1] \times [0, H-1]$. No formato planar semi-intercalado \texttt{NV21} do Android (Seção~\ref{sec:nv21}), os primeiros $W \times H$ bytes contêm o plano de luminância $Y$ em resolução total, seguidos por blocos intercalados de crominância $V$ e $U$ subamostrados em $2 \times 2$:
\begin{equation}
    i_Y = s_y\,W + s_x, \qquad
    i_{ VU } = W\,H + \lfloor s_y / 2 \rfloor\,W + 2\,\lfloor s_x / 2 \rfloor.
    \label{eq:indices_nv21}
\end{equation}
A partir de $Y = \text{\texttt{bytes}}[i_Y]$, $V = \text{\texttt{bytes}}[i_{VU}] - 128$ e $U = \text{\texttt{bytes}}[i_{VU} + 1] - 128$, converte-se a cor para RGB pela forma de faixa completa da BT.601 (Equação~\ref{eq:yuv_rgb_completa}), a mesma com que os JPEG da base MPIIGaze foram decodificados no treino, limita-se cada componente a $[0, 255]$, divide-se por 255 e escrevem-se as três componentes diretamente nas posições $i$, $S^2 + i$ e $2S^2 + i$ ($i = d_y S + d_x$) do \texttt{Float32List} de entrada do detector, já na ordem planar de canais (\texttt{NCHW}, Seção~\ref{sec:rgb}) exigida pelo modelo.

\subsection{Detecção da Pupila com YOLO26n} \label{sec:met_pupila}

A localização da pupila no recorte de $128 \times 128$ pixels é realizada pela variante \textit{nano} do YOLO26 (\texttt{YOLO26n}, Seção~\ref{sec:yolo})~\cite{yolo26}. O modelo foi treinado a partir dos pesos pré-treinados oficiais da Ultralytics sobre o subconjunto anotado (\textit{Annotation Subset}) da base MPIIGaze~\cite{zhang2015mpiigaze}, composto por 10.654 imagens capturadas por \textit{webcams} durante o uso diário de \textit{notebooks} por 15 participantes (\texttt{p00} a \texttt{p14}), sob iluminação ambiente não controlada. Como o MPIIGaze anota apenas o ponto central da pupila e os cantos dos olhos, sem máscaras de segmentação nem caixas delimitadoras, a preparação do conjunto de dados (\texttt{01\_\allowbreak prepare\_\allowbreak dataset.py}) aplica a transformação da Equação~\ref{eq:afim_recorte} aos dois olhos de cada imagem e sintetiza, em torno do centro anotado, uma caixa quadrada de lado igual a $20\%$ da largura aparente do olho (\texttt{pupilBoxRatio} $= 0{,}20$). Como a largura do olho se cancela com o fator de escala do recorte, toda caixa de referência possui lado constante igual a $0{,}20 / 1{,}6 = 0{,}125$ da imagem ($16 \times 16$ pixels no recorte de $128 \times 128$). Para evitar vazamento de características individuais entre treino e avaliação (Seção~\ref{sec:ml}), os 21.308 recortes foram divididos estritamente por participante: 11 participantes no treinamento (\texttt{p00}--\texttt{p10}, 16.780 recortes), 2 na validação (\texttt{p11}--\texttt{p12}, 2.196 recortes) e 2 reservados exclusivamente para teste (\texttt{p13}--\texttt{p14}, 2.332 recortes).

O treinamento (\texttt{02\_train.py}) foi configurado com classe única (\texttt{single\_cls}), lotes de 64 imagens, até 150 épocas e paciência de 25 épocas para parada antecipada (\textit{early stopping}). Durante o treinamento, aplicou-se aumento artificial de dados com translações aleatórias de até $15\%$ (\texttt{translate} $= 0{,}15$), variações de escala de até $30\%$ (\texttt{scale} $= 0{,}30$), rotações de $\pm 5^\circ$, espelhamento horizontal ($50\%$) e variações de brilho de até $60\%$ no espaço HSV (\texttt{hsv\_v} $= 0{,}60$). As perturbações de translação e escala simulam imprecisões nos cantos fornecidos pelo Face Mesh no primeiro estágio e impedem que a rede aprenda o atalho trivial de prever a pupila sempre no centro do recorte; em contrapartida, a técnica de mosaico foi desativada (\texttt{mosaic} $= 0{,}0$), pois reduziria cada quadrante a $64 \times 64$ pixels e degradaria um alvo de apenas 16 pixels de lado.

\begin{figure}[!htbp]
    \centering
    \includegraphics[width=\textwidth]{yolo_results.png}
    \caption{Evolução das funções de perda (\texttt{box\_loss}, \texttt{cls\_loss} e \texttt{l1\_loss}) nos conjuntos de treinamento e validação e das métricas de precisão, revocação, $\text{mAP@0,5}$ e $\text{mAP@0,5:0,95}$ ao longo das 38 épocas de treinamento do detector \texttt{YOLO26n}.}
    \label{fig:yolo_results}
\end{figure}

Como mostra a \refFig{fig:yolo_results}, as três componentes de perda no treinamento caíram rapidamente nas primeiras épocas e estabilizaram-se em torno de $1{,}45$ (\texttt{box\_loss}), $1{,}61$ (\texttt{cls\_loss}) e $0{,}011$ (\texttt{l1\_loss}). No conjunto de validação (participantes inéditos \texttt{p11} e \texttt{p12}), o modelo atingiu seu melhor desempenho na época~13 (\texttt{best.pt}), com as seguintes métricas (definidas na Seção~\ref{sec:metricas_visao}): precisão de $0{,}445$, revocação de $0{,}664$, $\text{mAP@0,5} = 0{,}496$ e $\text{mAP@0,5:0,95} = 0{,}162$; após essa época, as perdas de validação (\texttt{val/box\_loss} e \texttt{val/cls\_loss}) passaram a subir gradualmente enquanto as perdas de treino seguiam em queda, caracterizando o início de sobreajuste aos 11 indivíduos do conjunto de treinamento e acionando a parada antecipada na época~38, com restauração automática dos pesos da época~13.

\begin{figure}[!htbp]
    \centering
    \includegraphics[width=0.49\textwidth]{yolo_pr_curve.png}\hfill
    \includegraphics[width=0.49\textwidth]{yolo_f1_curve.png}
    \caption{Curvas de avaliação do modelo \texttt{YOLO26n} (\texttt{best.pt}) nos participantes de teste (\texttt{p13} e \texttt{p14}): à esquerda, a curva de precisão--revocação ($\text{mAP@0,5} = 0{,}601$); à direita, a curva de pontuação $F_1$ em função do limiar de confiança, com pico de $0{,}57$ no limiar $0{,}166$.}
    \label{fig:yolo_curvas_avaliacao}
\end{figure}

Na avaliação final sobre os 2.332 recortes dos participantes de teste \texttt{p13} e \texttt{p14} (\texttt{03\_\allowbreak evaluate.py}), o modelo alcançou $\text{mAP@0,5} = 0{,}601$ e $\text{mAP@0,5:0,95} = 0{,}210$ (\refFig{fig:yolo_curvas_avaliacao}, à esquerda). Como a caixa delimitadora é sintética e de proporção fixa em torno de um ponto, o mAP avalia apenas a sobreposição regional da caixa (Equação~\ref{eq:iou}), e não a exatidão da localização da pupila; a métrica relevante para o rastreamento ocular é a distância euclidiana entre o centro predito e o centro anotado no quadro original, normalizada pela largura do olho ($e_{\text{norm}} = e_{\text{recorte}}\,r / S$), conforme discutido na Seção~\ref{sec:metricas_visao}. Selecionando a caixa de maior confiança em cada recorte, o detector localizou a pupila em $100\%$ das 2.332 imagens de teste, com erro mediano de $4{,}46\%$ da largura do olho ($2{,}64$~pixels no quadro de $1280 \times 720$), média de $4{,}91\%$, percentil~90 de $8{,}53\%$ e percentil~95 de $9{,}86\%$, ficando abaixo de $5\%$ da largura do olho em $57{,}5\%$ das amostras e abaixo de $10\%$ em $95{,}2\%$ delas. O erro mediano manteve-se consistente entre os dois indivíduos inéditos ($4{,}56\%$ em \texttt{p13} e $4{,}38\%$ em \texttt{p14}), confirmando a generalização para novos rostos.

A curva de pontuação $F_1$ (Equação~\ref{eq:precisao_revocacao}) em função do limiar de confiança (\refFig{fig:yolo_curvas_avaliacao}, à direita) fundamenta uma decisão direta na implementação do aplicativo: como a fronteira da pupila sob luz visível é difusa e a caixa de treinamento é sintética, o modelo é sistematicamente subconfiante, concentrando suas pontuações em torno de $0{,}20$ (com pico de $F_1 = 0{,}57$ em $c = 0{,}166$ e valor máximo observado de $0{,}288$). A adoção do limiar padrão de detectores de objetos ($c = 0{,}25$) descartaria cerca de $97\%$ das detecções válidas. Como todo recorte entregue ao detector no aplicativo provém de um olho já localizado pelo Face Mesh e contém exatamente uma pupila, o papel do detector não é decidir se há uma pupila na imagem, mas apontar a melhor estimativa de sua posição. Por isso, o método \texttt{detect} da classe \texttt{PupilDetector} adota um limiar mínimo baixo (\texttt{defaultConfidence} $= 0{,}01$), percorre as 300 linhas candidatas do tensor de saída, seleciona a linha de maior confiança $c$ (\textit{top-1}) e devolve o ponto médio da caixa delimitadora:
\begin{equation}
    \mathbf{p}_{\text{recorte}} = \left(\frac{x_1 + x_2}{2},\, \frac{y_1 + y_2}{2}\right), \qquad
    \mathbf{p} = \mathbf{M}^{-1}\mathbf{p}_{\text{recorte}},
    \label{eq:centro_pupila}
\end{equation}
onde $\mathbf{p}$ já está expresso nas coordenadas do quadro bruto da câmera.

Para execução local no Android por meio do pacote \texttt{tflite\_flutter}, a rede treinada foi exportada para o formato LiteRT em precisão de ponto flutuante de 32~bits (\texttt{pupil\_128.tflite}, 9,7~MB), configurada com duas \textit{threads} de CPU. A exportação em FP32 foi escolhida porque preserva no próprio grafo estático as operações finais de decodificação livre de NMS do YOLO26, recebendo um tensor \texttt{[1, 3, 128, 128]} e entregando diretamente o tensor pós-processado \texttt{[1, 300, 6]}, em que cada uma das 300 linhas contém $(x_1, y_1, x_2, y_2, c, k)$, correspondentes aos cantos da caixa delimitadora em pixels do recorte, ao grau de confiança $c$ e à classe $k$. Experimentos no dispositivo mostraram que a quantização dinâmica híbrida (pesos em INT8 e ativações em FP32) falha na versão da biblioteca nativa embutida em \texttt{tflite\_flutter}, que não materializa tensores intermediários híbridos em tempo de execução, ao passo que a quantização inteira completa (INT8) omite o pós-processamento do grafo e devolveria a saída bruta \texttt{[1, 5, 336]}, exigindo implementar a decodificação e a supressão não máxima em Dart.

Na implementação de \texttt{PupilDetector}, constatou-se por medição no \textit{tablet} que a chamada nativa ao intérprete (\texttt{invoke}) consumia cerca de 15~ms por olho, mas a API de alto nível do \texttt{tflite\_flutter} baseada em listas aninhadas em Dart (\texttt{run} e \texttt{reshape}) acrescentava 53~ms por inferência apenas para copiar e converter estruturas de dados (\textit{marshalling}), elevando o custo total para 68,2~ms. Para eliminar essa sobrecarga, os objetos dos tensores de entrada e de saída são obtidos uma única vez no construtor e manipulados como blocos contíguos de memória: o método \texttt{detect} atribui os bytes do \texttt{inputBuffer} diretamente a \texttt{\_input.data} via visão \texttt{Uint8List}, aciona \texttt{\_interpreter.invoke()} e interpreta \texttt{\_output.data} no lugar como uma \texttt{Float32List} plana de $1.800$ elementos, reduzindo o tempo total dos dois olhos para 24,8~ms por quadro. No carregamento do modelo (\texttt{load}), executam-se seis inferências de aquecimento sobre um \textit{buffer} zerado para forçar a alocação imediata das arenas nativas do LiteRT e validar a integridade do arquivo antes que a câmera comece a enviar quadros.

\subsection{Vetor Pupila--Canto, Calibração e Regressão Polinomial} \label{sec:met_calibracao}

Com o centro da pupila $\mathbf{p}$ e os cantos suavizados $\mathbf{c}_A$ e $\mathbf{c}_B$ no mesmo sistema de coordenadas do quadro bruto, a função \texttt{eyeFeature} em \texttt{gaze\_mapper.dart} calcula para cada olho o vetor pupila--canto normalizado definido na Equação~\ref{eq:vetor_pupila_canto}. Quando as duas pupilas são localizadas no quadro (o que ocorre na quase totalidade dos quadros com rosto visível), o \texttt{GazeTracker} toma a média aritmética dos vetores do olho direito e do olho esquerdo:
\begin{equation}
    \mathbf{v} = (v_x, v_y) = \frac{\mathbf{v}_{\text{dir}} + \mathbf{v}_{\text{esq}}}{2},
    \label{eq:vetor_medio}
\end{equation}
reduzindo o desvio-padrão do ruído independente de detecção por um fator de $\sqrt{2}$ e mantendo o rastreamento funcional caso um dos olhos seja momentaneamente ocluído. O vetor resultante é armazenado no campo \texttt{lastFeature}, que alimenta tanto a tela de calibração quanto a predição em tempo real.

Para mapear $\mathbf{v}$ nas coordenadas normalizadas da tela $(x_{\text{tela}}, y_{\text{tela}}) \in [0, 1]^2$, é indispensável calibrar o sistema para o usuário e a posição atuais (Seção~\ref{sec:eyetracking}). A \texttt{Calibration\allowbreak Screen} (\refFig{fig:tela_calibragem}) apresenta sequencialmente uma grade regular de $5 \times 5$ alvos (25 posições) gerada por \texttt{calibrationTargets}, espaçados uniformemente entre $10\%$ e $90\%$ da largura e da altura da tela (\texttt{margin} $= 0{,}10$). Essa margem afasta os alvos das bordas extremas: testes no \textit{tablet} de 10,5~polegadas mostraram que alvos muito próximos às quinas levam o usuário a girar a cabeça em vez de rotacionar apenas o globo ocular, degradando as amostras periféricas, enquanto o comportamento suave do polinômio permite extrapolar até a borda com erro menor do que calibrar sobre pontos corrompidos.

Em cada alvo da grade, um temporizador leve (\texttt{\_Poller}, com passo de 33~ms) divide a permanência em duas fases:
\begin{enumerate}
    \item um intervalo inicial de estabilização de 700~ms (\texttt{\_settle}), durante o qual o anel de progresso permanece vazio e nenhuma amostra é guardada, dando tempo para a sacada ocular sair do alvo anterior e fixar o novo ponto; e
    \item uma janela de coleta de 1.200~ms (\texttt{\_collect}), na qual o anel âmbar se preenche e cada novo vetor \texttt{lastFeature} distinto do anterior é acrescentado à lista de amostras do alvo.
\end{enumerate}
Concluída a janela de coleta, se ao menos 5 amostras tiverem sido registradas, o ponto representativo do alvo é obtido pela mediana das coordenadas $v_x$ e $v_y$ das amostras (\texttt{\_medianOf}), protegendo o alvo contra piscadas ou detecções atípicas. A partir do sexto alvo válido coletado, a tela já reajusta um mapeamento provisório (\texttt{\_preview}) a cada novo alvo e desenha um círculo ciano na posição mediana das últimas 8 predições, oferecendo retorno visual imediato da evolução da calibração sem competir com o ponto branco que o usuário deve fixar.

Ao final dos 25 alvos (exigindo-se um mínimo de 6 alvos válidos para prosseguir), a função \texttt{fitGazeMap} ajusta os seis coeficientes por eixo do polinômio de segunda ordem da Equação~\ref{eq:regressao_polinomial}. Definindo para o $n$-ésimo alvo o vetor de características expandido $\mathbf{d}_n = (1,\, v_{x,n},\, v_{y,n},\, v_{x,n}^2,\, v_{y,n}^2,\, v_{x,n}v_{y,n})^\top$ e empilhando os $N$ alvos nas linhas da matriz de projeto $\mathbf{A}$ (de dimensão $N \times 6$), os vetores de seis coeficientes $\mathbf{a}$ e $\mathbf{b}$ dos eixos horizontal e vertical são obtidos pelas equações normais com regularização \textit{ridge} (Equação~\ref{eq:ridge}):
\begin{equation}
    \left(\mathbf{A}^\top \mathbf{A} + \lambda \mathbf{I}'\right)\mathbf{a} = \mathbf{A}^\top \mathbf{x}_{\text{alvo}}, \qquad
    \left(\mathbf{A}^\top \mathbf{A} + \lambda \mathbf{I}'\right)\mathbf{b} = \mathbf{A}^\top \mathbf{y}_{\text{alvo}},
    \label{eq:equacoes_normais_ridge}
\end{equation}
em que $\lambda = 10^{-4}$ (\texttt{kRidge}) estabiliza os termos quadráticos quando os pontos medidos apresentam quase-colinearidade e $\mathbf{I}' = \operatorname{diag}(0, 1, 1, 1, 1, 1)$ aplica a penalização apenas aos cinco coeficientes angulares e quadráticos, preservando o termo constante ($a_0, b_0$) para não deslocar artificialmente a estimativa em direção à origem da tela. Como o sistema linear de cada eixo tem dimensão fixa $6 \times 6$, ele é resolvido diretamente em Dart por eliminação de Gauss--Jordan com pivoteamento parcial (\texttt{\_solve}, Seção~\ref{sec:gauss}), sem dependência de bibliotecas externas de álgebra linear, lançando um \texttt{StateError} caso um pivô seja inferior a $10^{-12}$ (alvos degenerados).

Mesmo com a mediana por alvo e o afastamento das bordas, um único alvo em que o usuário tenha desviado o olhar distorce o ajuste de mínimos quadrados sobre todos os demais. Para tornar a calibração tolerante a falhas de fixação, a função \texttt{fitGazeMapRobust} executa um descarte iterativo de pontos discrepantes (\textit{outliers}): parte do ajuste com todos os $N$ alvos coletados, calcula a mediana dos erros euclidianos de reconstrução na tela e, por até \texttt{maxDrop} $= 2$ iterações (enquanto restarem ao menos 7 alvos), identifica o alvo de maior erro individual, reajusta o polinômio sem ele e conserva a remoção apenas se o resíduo mediano diminuir estritamente. Nos dados reais coletados no \textit{tablet} e fixados na suíte de testes (\texttt{gaze\_mapper\_test.dart}), esse descarte reduziu o resíduo mediano de calibração de $15{,}4\%$ para $12{,}1\%$ da tela.

\subsection{Filtragem de Kalman e Emissão de Comandos por Permanência} \label{sec:met_kalman_dwell}

Com o mapeamento calibrado atribuído ao \texttt{GazeTracker}, cada novo vetor $\mathbf{v}$ é projetado na tela por $\mathbf{z}_k = \text{\texttt{gazeMap.predict}}(\mathbf{v})$ e entregue ao filtro \texttt{GazeKalman} (\texttt{gaze\_kalman.dart}), que atenua o tremor quadro a quadro antes da decisão de comando. A classe mantém duas instâncias independentes de \texttt{\_AxisKalman}, uma para a coordenada horizontal e outra para a vertical, implementando de forma escalar as equações de predição e correção do modelo de velocidade constante (Equações~\ref{eq:kalman_pred_x} a~\ref{eq:kalman_matrizes}). No primeiro quadro após a calibração (ou após uma chamada a \texttt{reset}), o filtro de cada eixo é inicializado na própria medida inicial ($p_0 = z_0$), com velocidade nula ($\dot{p}_0 = 0$) e matriz de covariância diagonal $P_{00} = R$, $P_{01} = 0$ e $P_{11} = 1$. Nos quadros subsequentes, o intervalo $\Delta t$ decorrido desde a última atualização é calculado pelo relógio do sistema e limitado ao intervalo $[0{,}001\text{~s},\, 0{,}25\text{~s}]$, evitando que uma pausa momentânea no processamento produza uma extrapolação excessiva. Os hiperparâmetros adotados no código são a variância de medição $R = \sigma_z^2 = 0{,}075$ (\texttt{measurementVariance}) e a variância da aceleração $\sigma_a^2 = 1{,}0$ (\texttt{accelerationVariance}).

O ponto suavizado $(x, y) \in [0, 1]^2$ devolvido pelo filtro de Kalman é então classificado em uma das cinco ações de \texttt{GazeCommand} pela função \texttt{commandFor}. As quatro direções cardeais correspondem a faixas periféricas ao longo das bordas da tela, parametrizadas por \texttt{GazeBands} com largura horizontal $b_h = 0{,}33$ e altura vertical $b_v = 0{,}33$. Para resolver de forma contínua e simétrica os casos em que o olhar recai próximo às quinas (onde uma faixa horizontal e uma vertical se sobrepõem), calcula-se para cada borda a penetração normalizada do ponto dentro da respectiva faixa:
\begin{equation}
    \begin{aligned}
        s_{\text{\texttt{left}}} &= \frac{b_h - x}{b_h}, &\qquad
        s_{\text{\texttt{right}}} &= \frac{x - (1 - b_h)}{b_h}, \\[1ex]
        s_{\text{\texttt{up}}} &= \frac{b_v - y}{b_v}, &\qquad
        s_{\text{\texttt{down}}} &= \frac{y - (1 - b_v)}{b_v}.
    \end{aligned}
    \label{eq:penetrar_faixas}
\end{equation}
Se ao menos uma das quatro grandezas da Equação~\ref{eq:penetrar_faixas} for estritamente positiva, vence a direção de maior penetração relativa; caso todas sejam menores ou iguais a zero (isto é, quando $x \in [b_h, 1 - b_h]$ e $y \in [b_v, 1 - b_v]$, no terço central da tela), a função devolve \texttt{GazeCommand.select}.

Por fim, o método \texttt{\_updateDwell} converte a sequência de regiões observadas quadro a quadro em acionamentos discretos por tempo de permanência (\textit{dwell time}). O rastreador guarda o comando sob foco (\texttt{\_dwellCommand}) e o instante em que o olhar ingressou nessa região (\texttt{\_dwellStart}). Se o comando do quadro atual difere de \texttt{\_dwellCommand}, o foco mudou de região: atualiza-se \texttt{\_dwellCommand} e reinicia-se \texttt{\_dwellStart} para o instante presente. Se o comando permanece o mesmo e o tempo decorrido atinge o limiar configurado (\texttt{dwell} $= 1.500$~ms), o comando é emitido no fluxo \texttt{commands} --- que aciona \texttt{\_handleCommand} na \texttt{CommunicatorScreen} com origem \texttt{InputSource.gaze} --- e \texttt{\_dwellStart} é imediatamente rearmado para o instante atual, exigindo mais 1,5~s completo caso o usuário deseje repetir o mesmo comando. Em paralelo, todo quadro processado emite no fluxo \texttt{samples} um objeto \texttt{GazeSample} contendo o ponto suavizado, a região atual e os sinalizadores \texttt{faceFound} e \texttt{eyesFound}, a partir dos quais a tela principal realça o botão observado (\texttt{activeCommand}), desenha o cursor de depuração e atualiza o ícone de diagnóstico do cabeçalho.

  \section{Síntese de Voz (\texttt{tts/})} \label{sec:met_tts}

O pacote \texttt{tts/} isola a geração de fala offline descrita na Seção~\ref{sec:tts} por trás de uma interface comum:

\begin{itemize}
    \item \texttt{SpeechSynthesizer} e \texttt{SpeechSynthesizerState}: definem a interface abstrata do sintetizador (\texttt{initialize}, \texttt{speak}, \texttt{stop}, \texttt{dispose} e o estado observável \texttt{state}) e a exceção \texttt{SpeechSynthesisException}. Como a \texttt{CommunicatorScreen} depende apenas dessa abstração, os testes automatizados podem substituir o sintetizador real por uma implementação simulada em memória, sem carregar o modelo nem os \textit{plugins} nativos.
    \item \texttt{KokoroSynthesizer}: implementação que executa o modelo Kokoro (voz \texttt{pf\_dora}) com o pacote \texttt{sherpa\_onnx}. Copia os arquivos do modelo dos \textit{assets} para o diretório local do aplicativo na primeira execução (validando-os pelo \texttt{manifest.json}), mantém o sintetizador nativo em uma \textit{isolate} dedicada para não bloquear a interface nem o rastreamento ocular durante a geração do áudio, grava o resultado em um arquivo \texttt{WAV} temporário reproduzido via \texttt{audioplayers} e impede a sobreposição de falas simultâneas.
\end{itemize}

\begin{figure}[!htbp]
    \centering
    \includegraphics[width=\textwidth]{audio}
    \caption{Fluxo de síntese de voz a partir do acionamento da tecla FALAR, destacando a fonemização pelo \texttt{espeak-ng} e a inferência do modelo \texttt{Kokoro} executadas pelo \texttt{sherpa-onnx} em uma \textit{isolate} dedicada (azul) e a gravação e reprodução do arquivo \texttt{WAV} temporário (verde).}
    \label{fig:audio}
\end{figure}

Como mostra a \refFig{fig:audio}, quando o usuário confirma a tecla FALAR, o texto da mensagem percorre um pipeline local composto pela conversão fonética, pela inferência neural em segundo plano e pela reprodução do sinal de áudio resultante. As subseções a seguir detalham o contrato de abstração e o controle de estados na interface (Seção~\ref{sec:met_tts_contrato}), o empacotamento e a materialização verificada dos arquivos do modelo no dispositivo (Seção~\ref{sec:met_tts_modelo}) e a execução assíncrona em \textit{isolate} dedicada até a reprodução do áudio (Seção~\ref{sec:met_tts_isolate}).

\subsection{Contrato de Abstração e Controle de Estados} \label{sec:met_tts_contrato}

Definida em \texttt{tts/speech\_synthesizer.dart}, a interface abstrata \texttt{SpeechSynthesizer} desacopla a tela principal dos detalhes do motor nativo de síntese. O contrato expõe quatro métodos assíncronos --- \texttt{initialize}, \texttt{speak}, \texttt{stop} e \texttt{dispose} --- e um notificador observável (\texttt{ValueListenable<\allowbreak SpeechSynthesizerState>}), além da exceção \texttt{Speech\allowbreak Synthesis\allowbreak Exception} para encapsular falhas de preparação, inferência ou reprodução.

O ciclo de vida do sintetizador é representado pela enumeração \texttt{Speech\allowbreak Synthesizer\allowbreak State}, que assume seis estados possíveis:
\begin{itemize}
    \item \texttt{idle}: sintetizador disponível e aguardando uma nova frase;
    \item \texttt{preparing}: extração dos arquivos do modelo e inicialização do motor nativo na primeira chamada;
    \item \texttt{synthesizing}: conversão do texto em fonemas e execução da rede neural para gerar as amostras de áudio;
    \item \texttt{playing}: reprodução do arquivo de áudio gerado;
    \item \texttt{error}: falha recuperável ocorrida em qualquer etapa anterior; e
    \item \texttt{disposed}: sintetizador encerrado definitivamente junto com a tela.
\end{itemize}
Os três estados intermediários de operação (\texttt{preparing}, \texttt{synthesizing} e \texttt{playing}) são agrupados pela propriedade booleana \texttt{isBusy}. Na \texttt{CommunicatorScreen}, um ouvinte registrado sobre \texttt{synthesizer.state} reconstrói o teclado sempre que o estado muda, atualizando o rótulo da primeira tecla de ação para refletir a etapa em curso (\texttt{CARREGANDO}, \texttt{GERANDO}, \texttt{FALANDO} ou \texttt{FALAR}) e desabilitando a tecla enquanto \texttt{isBusy} for verdadeiro ou a mensagem estiver vazia.

Essa sinalização resolve um problema próprio de interfaces acionadas por tempo de permanência: após disparar a tecla FALAR, o usuário frequentemente mantém o olhar sobre a mesma região por mais alguns instantes enquanto aguarda o início da voz. Sem o bloqueio por estado ocupado --- reforçado no próprio método \texttt{speak} do \texttt{KokoroSynthesizer}, que rejeita chamadas concorrentes enquanto \texttt{\_activeSpeak} não for nulo ---, um segundo ciclo de permanência (\textit{dwell}) dispararia uma nova síntese sobreposta à primeira. Além disso, como a tela depende apenas de \texttt{SpeechSynthesizer}, a suíte de testes de interface (\texttt{widget\_test.dart}) injeta uma implementação leve em memória (\texttt{\_Fake\allowbreak Speech\allowbreak Synthesizer}), permitindo verificar a preservação da mensagem após a fala e a mudança dos rótulos da tecla FALAR sem carregar os 355~MB do modelo nem invocar bibliotecas nativas.

\subsection{Preparação Offline e Materialização Verificada do Modelo} \label{sec:met_tts_modelo}

Para manter o repositório Git leve e, ao mesmo tempo, garantir reprodutibilidade e funcionamento integralmente offline no \textit{tablet}, a preparação dos recursos de voz é automatizada pelo \textit{script} \texttt{tool/fetch\_kokoro\_model.py}. O \textit{script} baixa o pacote oficial \texttt{kokoro-\allowbreak multi-\allowbreak lang-\allowbreak v1\_0.tar.bz2} (349,9~MB) distribuído pelo projeto \texttt{sherpa-onnx}~\cite{sherpa_onnx}, valida o tamanho e o resumo criptográfico SHA-256 do arquivo compactado e extrai para \texttt{assets/\allowbreak models/\allowbreak kokoro\_pt/} apenas os dez arquivos necessários à síntese em português brasileiro:
\begin{itemize}
    \item \texttt{model.onnx} (325,6~MB): grafo estático em ponto flutuante de 32~bits (FP32) do modelo Kokoro~\cite{kokoro82m};
    \item \texttt{voices.bin} (28,2~MB) e \texttt{tokens.txt} (687~bytes): tabela binária dos vetores de estilo de voz e mapeamento dos símbolos fonéticos para identificadores numéricos; e
    \item sete arquivos de dados do \texttt{espeak-ng} (\texttt{phondata}, \texttt{phonindex}, \texttt{phontab}, \texttt{intonations\allowbreak}, \texttt{pt\_dict}, \texttt{lang/roa/pt} e \texttt{lang/roa/pt-BR}): subconjunto de 715~KB que substitui os 355 arquivos (19~MB) do diretório \texttt{espeak-ng-data} completo, produzindo saída fonética idêntica para o idioma \texttt{pt-br}.
\end{itemize}
Após validar o tamanho e o resumo SHA-256 de cada arquivo extraído, o \textit{script} grava o arquivo \texttt{manifest.json} com os metadados de versão e integridade (que também podem ser auditados sem acesso à rede com a opção \texttt{--verify-only}).

Em dispositivos Android, os arquivos declarados em \texttt{assets/} ficam empacotados dentro do arquivo \texttt{APK} e não possuem caminhos diretos no sistema de arquivos do sistema operacional, requisito exigido pela biblioteca nativa em C++ do \texttt{sherpa-onnx} carregada via \texttt{dart:ffi}. Por essa razão, na primeira chamada a \texttt{initialize}, o método \texttt{\_materializeModel} do \texttt{KokoroSynthesizer} copia os recursos do pacote para o diretório privado de suporte do aplicativo (\texttt{getApplicationSupportDirectory/kokoro\_pt}). Para não repetir a cópia de cerca de 354~MB a cada abertura do aplicativo, a função \texttt{\_cacheIsCurrent} compara o \texttt{manifest.json} gravado no armazenamento local com o manifesto embarcado nos \textit{assets} e verifica se todos os dez arquivos listados existem no diretório local com o tamanho exato em bytes especificado no manifesto. Quando a cópia é necessária, cada arquivo é gravado de forma atômica (\texttt{\_writeAtomically}) em um arquivo temporário com sufixo \texttt{.part} (com descarga forçada em disco via \texttt{flush: true}) e renomeado para o nome definitivo apenas após a conclusão, sendo o \texttt{manifest.json} local gravado por último. Dessa forma, caso o aplicativo seja fechado ou perca energia no meio da primeira inicialização, a próxima execução detecta a ausência ou divergência do manifesto e refaz a materialização sem risco de alimentar o motor nativo com arquivos truncados.

\subsection{Execução em \textit{Isolate} Dedicada e Reprodução de Áudio} \label{sec:met_tts_isolate}

A síntese neural de uma frase em CPU leva da ordem de centenas de milissegundos a alguns segundos, a depender do comprimento do texto. Se a chamada síncrona da biblioteca nativa (\texttt{tts.generate}) fosse executada na \textit{isolate} principal do Flutter, o laço de eventos da interface gráfica e o processamento dos quadros da câmera pelo \texttt{GazeTracker} ficariam congelados durante toda a geração do áudio. Para isolar esse custo, o método \texttt{\_spawnWorker} cria, durante a inicialização, uma \textit{isolate} trabalhadora persistente (\texttt{\_kokoroWorker}, identificada como \texttt{kokoro-tts}), que permanece ativa durante toda a vida útil da tela. Manter uma única \textit{isolate} de longa duração --- em vez de criar uma nova a cada frase --- evita o custo de recarregar em memória os 325,6~MB do grafo \texttt{model.onnx} a cada acionamento da tecla FALAR.

Ao iniciar, a \textit{isolate} trabalhadora carrega as ligações FFI (\texttt{sherpa.initBindings()}) e instancia o objeto nativo \texttt{OfflineTts} configurado com os caminhos locais materializados, idioma \texttt{pt-br}, provedor de execução \texttt{cpu} e duas \textit{threads} de processamento (\texttt{numThreads: 2}), número escolhido para dividir os núcleos do processador do \textit{tablet} sem competir de forma agressiva com o rastreador ocular. Em seguida, abre uma porta de recepção (\texttt{ReceivePort}) e envia o respectivo \texttt{SendPort} de comando de volta à \textit{isolate} principal em uma mensagem \texttt{ready}. A comunicação a partir desse ponto segue um protocolo assíncrono de requisição e resposta: a \textit{isolate} principal associa a cada pedido um identificador inteiro sequencial (\texttt{id}) e um \texttt{Completer} armazenado no mapa \texttt{\_requests}, enviando comandos tipados (\texttt{synthesize} ou \texttt{dispose}) pela porta da trabalhadora. Para garantir que falhas nativas nunca deixem promessas pendentes indefinidamente na interface, o \texttt{Isolate.spawn} registra portas de escuta para erros não tratados e encerramento abrupto (\texttt{onError} e \texttt{onExit}), que acionam \texttt{\_failWorker} para completar imediatamente todas as requisições em aberto com uma \texttt{SpeechSynthesisException} e permitir que uma chamada futura a \texttt{speak} reinicialize a \textit{isolate}.

Quando o usuário aciona a síntese de uma frase por \texttt{speak(text)}, o método executa a sequência completa de ponta a ponta:
\begin{enumerate}
    \item remove espaços excedentes da cadeia de entrada, rejeita textos vazios e garante que a inicialização do modelo e da \textit{isolate} trabalhadora esteja concluída;
    \item transita o estado para \texttt{synthesizing} e envia à \textit{isolate} trabalhadora o comando \texttt{synthesize} com o texto e o caminho de saída no subdiretório temporário \texttt{tts\_\allowbreak audio/} (\texttt{speech\_\allowbreak <timestamp>.wav});
    \item na \textit{isolate} trabalhadora, invoca \texttt{tts.generate} selecionando o identificador de locutor \texttt{\_voiceId} $= 42$ (correspondente à voz feminina brasileira \texttt{pf\_dora} na tabela do Kokoro~v1.0), verifica que o vetor de amostras não é vazio e grava o sinal PCM de 24~kHz no arquivo indicado por meio de \texttt{sherpa.writeWave}, devolvendo o caminho confirmado à \textit{isolate} principal;
    \item ao receber a resposta na \textit{isolate} principal (e caso \texttt{stop} não tenha sido solicitado durante a espera), muda o estado para \texttt{playing} e reproduz o arquivo local com o pacote \texttt{audioplayers} (\texttt{DeviceFileSource} com tipo MIME \texttt{audio/wav} e modo de liberação \texttt{ReleaseMode.stop}), aguardando o evento \texttt{onPlayerComplete}; e
    \item em um bloco \texttt{finally}, apaga imediatamente o arquivo \texttt{WAV} temporário do armazenamento e retorna o estado para \texttt{idle}.
\end{enumerate}
Como salvaguarda complementar, a inicialização também executa \texttt{\_removeStaleAudio} para remover eventuais arquivos \texttt{.wav} remanescentes de sessões anteriores encerradas de forma abrupta, e o método \texttt{dispose} interrompe qualquer reprodução ativa, envia o comando de encerramento à \textit{isolate} trabalhadora (com tempo limite de 5~s antes do encerramento forçado via \texttt{Isolate.immediate}) e libera a memória nativa do sintetizador (\texttt{tts.free()}).

  \section{Registro de Métricas e Avaliação (\texttt{metrics/})} \label{sec:met_metricas}

O pacote \texttt{metrics/} concentra o registro das sessões de uso e o cálculo das métricas de avaliação fundamentadas na Seção~\ref{sec:metricas} e detalhadas nas subseções a seguir:

\begin{itemize}
    \item \texttt{TypingEvent}: representa cada evento bruto ocorrido durante a digitação --- comando emitido (com indicação de origem pelo olhar ou por toque), tecla acionada, lista de sugestões exibida ou sugestão aceita --- marcado com o tempo decorrido desde o início da sessão e serializável em JSON.
    \item \texttt{TypingMetricsRecorder}: registra os eventos em tempo real durante o uso da \texttt{CommunicatorScreen} e salva o arquivo JSON da sessão no armazenamento do dispositivo a cada frase falada e ao encerrar a tela, incluindo tanto a sequência bruta de eventos quanto as métricas consolidadas.
    \item \texttt{TypingMetrics}, \texttt{PhraseMetrics} e \texttt{TypingCounts}: reconstroem, a partir da lista de \texttt{TypingEvent}, as frases delimitadas pelo acionamento da tecla FALAR e calculam para cada frase e para o conjunto da sessão as medidas de velocidade (WPM e tempo por comando), esforço (KSPC e economia de comandos em relação à digitação sem preenchimento automático, cuja linha de base é calculada por busca em largura segundo a Equação~\ref{eq:custo_base}), erros corrigidos e desempenho das sugestões (Seções~\ref{sec:met_velocidade} a~\ref{sec:met_agregacao}).
    \item \texttt{SimulatedTypist}: simula a digitação das frases do \textit{corpus} seguindo as mesmas regras de navegação e aceitação de sugestões da tela principal, com suporte à validação cruzada (\texttt{crossValidate}), à avaliação nas frases de treino (\texttt{evaluateInSample}) e à injeção de erros de digitação, conforme descrito na Seção~\ref{sec:met_avaliacao}.
\end{itemize}

Diferentemente da avaliação offline, que simula a digitação de frases conhecidas (Seção~\ref{sec:met_avaliacao}), o registro de métricas acompanha o uso real do aplicativo. Ele foi projetado com dois requisitos: não interferir na comunicação e permitir a análise do uso livre, em que não se sabe de antemão o que a pessoa quer escrever. As subseções a seguir descrevem o registro dos eventos (Seção~\ref{sec:met_eventos}), a unidade de análise (Seção~\ref{sec:met_unidade}), as medidas de velocidade e esforço (Seção~\ref{sec:met_velocidade}), a linha de base e a economia de comandos (Seção~\ref{sec:met_economia}), os erros corrigidos (Seção~\ref{sec:met_erros}), o desempenho das sugestões (Seção~\ref{sec:met_desempenho_sugestoes}) e a agregação dos resultados de uma sessão (Seção~\ref{sec:met_agregacao}).

\subsection{Registro de Eventos} \label{sec:met_eventos}

Cada abertura da tela principal inicia uma sessão de registro. A partir desse momento, a \texttt{CommunicatorScreen} informa ao \texttt{TypingMetricsRecorder} cada ação relevante, que é armazenada como um evento junto com o tempo decorrido desde o início da sessão, em milissegundos. Há quatro tipos de evento, resumidos na Tabela~\ref{tab:eventos}.

\begin{table}[htbp]
    \centering
    \caption{Tipos de evento registrados durante a digitação.}
    \label{tab:eventos}
    \begin{tabular}{lp{5cm}p{6cm}}
        \hline
        Tipo & Quando é registrado & Dados \\
        \hline
        \texttt{command} & a cada comando direcional ou de seleção & o comando e sua origem (olhar ou toque) \\
        \texttt{key} & a cada tecla acionada & a ação, o rótulo e a posição da tecla e a mensagem após o acionamento \\
        \texttt{suggestions} & sempre que as sugestões exibidas mudam & a mensagem e as sugestões \\
        \texttt{acceptance} & a cada sugestão inserida & a sugestão, sua posição na barra e a mensagem após a inserção \\
        \hline
    \end{tabular}
\end{table}

A origem de cada comando distingue os que vieram do rastreamento ocular dos que vieram da tela sensível ao toque, seja pelos botões direcionais exibidos na tela, que emitem os mesmos comandos do olhar, seja por um toque direto em uma tecla ou sugestão. O toque direto é registrado como um comando de seleção sem os movimentos que o precederiam, o que tem consequências para a economia de comandos (Seção~\ref{sec:met_economia}). As sugestões só são registradas quando mudam, e não a cada atualização da tela, para que o log não cresça sem necessidade.

O registro guarda apenas os eventos brutos, e todas as métricas são calculadas a partir deles pela classe \texttt{TypingMetrics}. Assim, se a definição de uma métrica mudar, as sessões já gravadas podem ser recalculadas sem que os testes sejam repetidos. A cada acionamento da tecla FALAR e ao fechamento da tela, a sessão inteira é gravada em um arquivo JSON, que sobrescreve a versão anterior e contém os eventos e as métricas calculadas até aquele momento. O arquivo recebe o nome \texttt{session\_AAAAMMDD\_HHMMSS.json}, com a data e a hora de início da sessão, e fica na pasta externa do aplicativo, de onde pode ser copiado para um computador pelo \textit{Android Debug Bridge} (\texttt{adb}) sem acesso de superusuário. Gravar a cada frase falada limita a perda, caso o aplicativo seja encerrado de forma abrupta, à frase em andamento. Uma falha na gravação apenas gera um aviso no log de depuração, de modo que o registro de métricas nunca interrompe a comunicação.

\subsection{Unidade de Análise: a Frase Falada} \label{sec:met_unidade}

As medidas da Seção~\ref{sec:metricas} são definidas sobre um texto transcrito $T$ e o tempo gasto para produzi-lo, mas no uso livre não há uma frase-alvo que delimite esse texto. A unidade de análise adotada é a frase falada: tudo o que acontece entre dois acionamentos da tecla FALAR. A decisão de falar marca o fim de uma mensagem e é o único momento em que o aplicativo sabe que o texto está pronto.

Como a mensagem permanece na tela depois de falada, o texto de uma frase é apenas o que foi produzido depois do FALAR anterior. Três regras completam a definição:
\begin{itemize}
    \item se a pessoa continua a mensagem anterior, só a parte nova é atribuída à frase: após falar ``EU QUERO ÁGUA'' e acrescentar ``GELADA'', a segunda frase é ``GELADA'';
    \item apagar o texto herdado da frase anterior, com APAGAR ou LIMPAR, é o começo de uma frase nova, e não a correção de um erro; e
    \item acionar FALAR de novo sem alterar a mensagem não cria uma frase.
\end{itemize}
O tempo da frase vai do primeiro comando depois do FALAR anterior até o FALAR que a encerra. Não são contados o intervalo entre frases, em que a pessoa pode estar ouvindo a fala sintetizada, descansando ou conversando, nem o tempo antes do primeiro comando da sessão, que inclui a calibração.

\subsection{Velocidade e Esforço} \label{sec:met_velocidade}

Sendo $|T|$ o número de caracteres da frase, incluindo os espaços entre as palavras, e $S$ o seu tempo, em segundos, a velocidade é calculada como
\begin{equation}
    \mathrm{WPM} = \frac{|T|}{S} \times 60 \times \frac{1}{5}.
    \label{eq:wpm_app}
\end{equation}
A diferença em relação à Equação~\ref{eq:wpm} é a ausência da subtração de um caractere. Na medição tradicional, o tempo começa no primeiro caractere, cuja entrada não é cronometrada~\cite{wobbrock2007measures}; aqui, ele começa no primeiro comando, de modo que a navegação até o primeiro caractere e a sua seleção já fazem parte de $S$. Também é registrado o tempo médio por comando, $S/C$, em que $C$ é o número de comandos da frase, o que permite comparar o ritmo da pessoa com o tempo de permanência de 1,5~s necessário para emitir cada comando pelo olhar.

O esforço é medido pelo KSPC de desempenho (Equação~\ref{eq:kspc}). Seguindo a interpretação ampla de ``toque''~\cite{wobbrock2007measures}, cada comando conta como uma ação, e a sequência de entrada $IS$ reúne os movimentos e as seleções da frase, vindos do olhar ou do toque, de modo que $|IS| = C$. Esse valor inclui a navegação: digitar uma letra custa, além da seleção, os comandos necessários para levar o foco até ela. Digitar ``SIM'' letra a letra a partir da tecla Q, por exemplo, exige 19 comandos, dos quais 4 são seleções (as três letras e FALAR), e o KSPC é de $19/3 \approx 6{,}33$ mesmo sem nenhum erro. Para permitir a comparação com a literatura, em que cada toque escolhe diretamente uma tecla, é calculada também uma versão que conta apenas as seleções, que nesse exemplo é de $4/3 \approx 1{,}33$. A avaliação offline (Seção~\ref{sec:met_avaliacao}), por sua vez, corresponde ao KSPC característico, pois mede o custo da técnica sobre um conjunto de frases, sem usuários.

\subsection{Linha de Base e Economia de Comandos} \label{sec:met_economia}

Para cada frase, a linha de base é o número mínimo de comandos para produzir o mesmo texto letra a letra, sem as sugestões, calculado pela Equação~\ref{eq:custo_base} com as distâncias obtidas por busca em largura (Seção~\ref{sec:bfs}). A tecla inicial $k_0$ é a última tecla acionada antes da frase, ou a tecla Q, em que o foco começa ao abrir a tela, e a sequência $k_1, \ldots, k_m$ é formada pelas letras do texto, com as palavras separadas pela tecla ESPAÇO, e termina na tecla FALAR. Duas adaptações são necessárias. Primeiro, o texto é normalizado como no modelo de linguagem (Seção~\ref{sec:met_normalizacao}), pois o teclado não tem acentos nem hífen: uma palavra como ``ÁGUA'', inserida por sugestão, custa na linha de base as letras A, G, U e A. Segundo, se o texto herdado da frase anterior foi apagado, a linha de base inclui um acionamento da tecla LIMPAR antes da primeira letra. Da mesma forma, a linha de base em seleções é o número de teclas da sequência.

A economia de comandos é então
\begin{equation}
    E = 1 - \frac{C}{C_{\text{base}}},
    \label{eq:economia}
\end{equation}
e a economia de seleções é calculada da mesma forma, com as seleções no lugar dos comandos. A frase ``EU QUERO BEBER ÁGUA'', partindo da tecla Q, pode ser escrita com 18 comandos: o E é digitado e as quatro palavras são inseridas pela barra de sugestões, a primeira completando o E. Letra a letra, a mesma frase exigiria 97 comandos, e a economia é de $1 - 18/97 \approx 81{,}4\%$. Em seleções, são 6 contra 20, uma economia de 70\%. Com 19 caracteres, o KSPC dessa frase é de $18/19 \approx 0{,}95$, menos de um comando por caractere.

Ao contrário da economia de toques da Seção~\ref{sec:metricas}, que compara o custo ideal de duas técnicas, a Equação~\ref{eq:economia} compara o que a pessoa de fato fez com o custo ideal da digitação sem sugestões. Por isso, ela pode ser negativa: comandos desperdiçados por imprecisão do rastreamento e erros corrigidos aumentam $C$ sem alterar $C_{\text{base}}$, que só depende do texto final. A economia junta, assim, em um único número, o ganho das sugestões e o custo da navegação imprecisa e dos erros, que só podem ser separados examinando o log de eventos.

Um toque direto em uma tecla é registrado como uma única seleção, sem os movimentos que a linha de base supõe. Em sessões digitadas dessa forma, a economia de comandos e o KSPC perdem o significado, e apenas as medidas baseadas em seleções permanecem comparáveis. Por isso, os testes devem ser feitos pelo olhar ou pelos botões direcionais da tela.

\subsection{Erros Corrigidos} \label{sec:met_erros}

Sem frase-alvo, não é possível saber se uma palavra que ficou na frase era a pretendida, e os erros não corrigidos da Equação~\ref{eq:taxas_erro} não podem ser medidos pelo aplicativo. Eles exigem testes com frases-alvo, em que a frase falada é comparada com a esperada, fora do aplicativo. Os erros corrigidos, por outro lado, ficam registrados nos eventos: cada caractere produzido na frase e depois apagado, com APAGAR ou LIMPAR, corresponde a um toque incorreto mas corrigido ($IF$). Como o texto final reúne os toques corretos e os incorretos não corrigidos, o denominador da Equação~\ref{eq:taxas_erro} equivale a $|T| + IF$, e a taxa de erros corrigidos é
\begin{equation}
    \text{taxa de erros corrigidos} = \frac{IF}{|T| + IF}.
    \label{eq:erros_corrigidos}
\end{equation}
Se, ao escrever ``SIM'', a pessoa aciona J em vez de M, apaga a letra e digita M, a frase tem 3 caracteres e 1 caractere apagado, e a taxa é de $1/(3+1) = 0{,}25$. Os toques de correção ($F$) são contados à parte, como o número de acionamentos de APAGAR e de LIMPAR.

Dois casos não contam como erro. O primeiro é apagar o texto herdado da frase anterior, que inicia uma frase nova (Seção~\ref{sec:met_unidade}). O segundo é a substituição da palavra parcial por uma sugestão, como ao aceitar ``CABEÇA'' depois de ``CSB'' (Seção~\ref{sec:met_tolerancia}): a sugestão reescreve a palavra, inclusive para inserir acentos, e tratar essa troca como erro faria toda palavra aceita parecer uma correção. Com isso, um erro de digitação consertado pela tolerância a erros não aparece na taxa, embora continue visível no log de eventos. São registradas, ainda, as sugestões rejeitadas, isto é, aceitas mas cuja palavra não permaneceu na frase falada, o que indica uma sugestão inserida por engano e depois apagada.

\subsection{Desempenho das Sugestões} \label{sec:met_desempenho_sugestoes}

Os eventos de sugestões exibidas permitem saber se e quando cada palavra da frase falada foi prevista. Uma palavra é considerada prevista depois de $j$ letras se, em algum momento em que havia $j$ letras digitadas na sua posição, ela estava entre as sugestões; $j = 0$ indica que ela foi sugerida como próxima palavra, antes de qualquer letra. As letras contadas incluem as erradas: como o modelo tolera erros de digitação, ``BEBER'' sugerida após ``BWB'' é uma palavra prevista depois de 3 letras. A partir disso, são calculadas:
\begin{itemize}
    \item a acurácia top-3 e a top-1 (Seção~\ref{sec:metricas}), frações das palavras que apareceram entre as três sugestões e como primeira sugestão, respectivamente, antes de serem concluídas;
    \item o número médio de letras digitadas até a palavra aparecer, entre as palavras previstas;
    \item a taxa de conclusão, fração das palavras inseridas por uma sugestão;
    \item a utilização, fração das palavras previstas que foram de fato aceitas, que indica se a pessoa aproveita as sugestões corretas quando elas aparecem; e
    \item a distribuição das posições, na barra, das sugestões aceitas.
\end{itemize}
Uma palavra que só apareceu nas sugestões depois de digitada por inteiro e foi então aceita não conta como inserida por sugestão. Nesse caso, a sugestão serve apenas de atalho para o espaço ou para o acento (``AGUA'' $\to$ ``ÁGUA''), sem poupar nenhuma letra. Na frase ``EU QUERO BEBER ÁGUA'', ``EU'' apareceu na segunda posição depois de 1 letra, e as outras três palavras, na primeira posição, antes de qualquer letra: a acurácia top-3 é de 1,0, a top-1 é de 0,75, a média de letras até a previsão é de 0,25 e a taxa de conclusão é de 1,0.

\subsection{Agregação da Sessão} \label{sec:met_agregacao}

Além das métricas de cada frase, o arquivo da sessão traz um resumo de todas elas. O resumo soma as contagens das frases (caracteres, comandos, linha de base, caracteres apagados, palavras previstas e assim por diante) e só então calcula as razões sobre os totais. Essa micro-média, em vez da média das razões de cada frase, dá o mesmo peso a cada caractere e a cada comando, de modo que uma frase curta, como ``SIM'', pesa no resultado da sessão na proporção do seu tamanho. Uma razão sem denominador, como a utilização em uma frase sem palavras previstas, é registrada como nula, e não como zero, para não ser confundida com um valor medido. O resumo informa, ainda, quantos comandos vieram do olhar e quantos do toque, o que permite verificar se a sessão foi de fato realizada no modo pretendido.
